% Navier--Stokes: costura em campo --- leitura pela Lei~4
% Tiffany — papers/navier-stokes.tex
% Compila com pdflatex navier-stokes.tex
\documentclass[11pt,a4paper]{article}
\usepackage[utf8]{inputenc}
\usepackage[T1]{fontenc}
\usepackage[portuguese]{babel}
\usepackage{amsmath,amssymb,amsthm}
\usepackage[margin=2.3cm]{geometry}
\usepackage{booktabs}
\usepackage[hidelinks]{hyperref}
\usepackage{xcolor}

\newcommand\code[1]{\texttt{\small #1}}
\newcommand\medido[1]{%
  \par\medskip\noindent{\small\color{gray}\textbf{Medido.} #1}\par\medskip}
\newtheorem{teorema}{Teorema}
\newtheorem{definicao}[teorema]{Definição}
\newtheorem{proposicao}[teorema]{Proposição}
\newcommand\dg{^{\dagger}}

\title{Navier--Stokes:\\a costura em campo --- leitura pela Lei~4}
\author{broca-so / Tiffany}
\date{}
\begin{document}
\maketitle

\begin{abstract}
Navier--Stokes é lido nesta arquitetura através da
\textbf{Lei~4} ($T+T^{\dg}$, $|\det|=1$) no palco de campo:
transporte (esquilo) + dissipação (gato) = costura.
O contínuo 3D é declarado \emph{fronteira}, não teorema.
Incluímos o resultado anunciado pela OpenAI (set.~2026) como
\textbf{claim} documentado, sem transformá-lo em validação
da arquitetura.
\end{abstract}

%==========================================================================
\section{Introdução}\label{sec:intro}
%==========================================================================

As equações de Navier--Stokes incompressíveis:
\begin{equation}\label{eq:NS}
\partial_t u+(u\cdot\nabla)u
=
\nu\Delta u-\nabla p+f,
\qquad
\nabla\cdot u=0,
\end{equation}
com $u:\mathbb R^3\times[0,T)\to\mathbb R^3$,
$\nu>0$, força $f$ suave compactamente apoiada quando relevante.

A pergunta clássica do Clay: para dados iniciais suaves,
existe $T=\infty$ com $u$ suave? Ou há \emph{blowup}
(singularidade em tempo finito)?

Este paper apresenta quatro camadas distintas:

\begin{description}
\item[A.~Matemática estabelecida] equação NS, incompressibilidade,
  projeção de Leray $\mathbb P$, operador de Stokes,
  identidade de energia padrão;
\item[B.~Resultado medido no projeto] identidades testadas em
  \code{tests/navier_corpo.js},
  \code{tests/histerese_navier.js},
  \code{tests/batuta_continuo.js},
  \code{tests/curvatura.c};
\item[C.~Estrutura canônica do projeto] Lei~4, $T+T^{\dg}$,
  $|\det|=1$, Transformada $\mathcal{F}$, gato/esquilo,
  costura, fronteira, histerese;
\item[D.~Interpretação] leitura da Lei~4 sobre a equação de campo,
  decomposição conservativo/dissipativo, leitura espectral,
  ponte com Yang--Mills --- \textbf{tudo explicitamente marcado
  como interpretação}.
\end{description}

%==========================================================================
\section{Lei~4 --- $T+T^{\dg}$ e $|\det|=1$}\label{sec:lei4}
%==========================================================================

A Lei~4 é a \textbf{tetral} ($n=4$) do Corpo Universal~U:
\begin{equation}\label{eq:lei4}
T+T^{\dg},\qquad |\det|=1.
\end{equation}
O que cada termo significa na arquitetura canónica:

\begin{description}
\item[$T+T^{\dg}$] operador da tetral; também escrito
  $T+T^{*}$; é a soma do operador com o seu dual.
  No contexto do dragão, $T+T^{*}=1$ (identidade).
\item[$|\det|=1$] métrica canónica de U: conservação de volume.
  É $\mathrm{I}_{11}$ (morfismos/estado), não $\mathrm{I}_{4}$
  (GLH).
\end{description}

\textbf{Não} é correto escrever ``$T+T^{\dg}$ é a estrela'' nem
``$|\det|=1$ é a estrela''. As estrelas são três autoridades
distintas (Obs.~($\star$three-canonicos)):
\begin{equation}\label{eq:tres-estrelas}
\operatorname{Star}(\mathcal{U})=\mathcal{D}
\quad\neq\quad
\operatorname{Star}(K)=\{x:x^{2}=x+1\}
\quad\neq\quad
\text{factor/métrica local}.
\end{equation}
Star(U)=D é a troca $(\oplus,\otimes)\mapsto(\otimes,\oplus)$
(bit $a=1$); $T+T^{\dg}$ é o operador da tetral;
$|\det|=1$ é a métrica de volume. Não se confundem.

A correspondência com Navier--Stokes é \textbf{leitura interna}
do projeto: não é, por si só, equivalência matemática
estabelecida com a formulação do Clay.

%==========================================================================
\section{Navier--Stokes padrão}\label{sec:ns-padrao}
%==========================================================================

Na formulação com operador de Stokes, define-se
\begin{equation}\label{eq:stokes}
A = -\mathbb P\Delta,
\end{equation}
$\mathbb P$ a projeção de Leray. Com esta convenção,
$A$ é \textbf{positivo} e o semigrupo do calor é
$e^{-\nu tA}$.

A equação fica
\begin{equation}\label{eq:NS-stokes}
\partial_t u + \nu A u = -\mathbb P((u\cdot\nabla)u) + \mathbb P f,
\end{equation}
e a forma energética (caso $f=0$) é
\begin{equation}\label{eq:energia-padrao}
\frac{d}{dt}\frac12\|u\|_2^2 = -\nu\langle Au,u\rangle
= -\nu\|\nabla u\|_2^2 \le 0.
\end{equation}

\textbf{Não} se escreve que o espectro de $A$ é negativo:
se $-A$ tem espectro negativo, $A$ tem espectro \textbf{positivo}.
O operador $-\nu A$ é que tem espectro negativo; não $A$.

Para o caso forçado:
\begin{equation}\label{eq:energia-forcada}
\frac{d}{dt}\frac12\|u\|_2^2
=
-\nu\|\nabla u\|_2^2 + \langle f,u\rangle,
\end{equation}
e a energia \textbf{não decresce estritamente} em geral
(o termo $\langle f,u\rangle$ pode injectar energia).
Isto não é prova de regularidade 3D.

%==========================================================================
\section{Gato / esquilo}\label{sec:gato-esquilo}
%==========================================================================

Os nomes \textbf{gato} e \textbf{esquilo} são vocabulário
estrutural do projeto, não nomenclatura padrão da literatura
de Navier--Stokes.
\begin{description}
\item[gato] $= $ convolução $\zeta$ (operador que sobe; conservativo)
\item[esquilo] $= $ deconvolução $\mu=\zeta^{-1}$ (operador que desce; dissipativo)
\end{description}
Conforme as definições canónicas encontradas em
\code{fisica.tex}~\S f{is:def:gato}.

Na leitura do projeto:
\begin{equation}\label{eq:costura-leitura}
\underbrace{\text{transporte/concentração}}_{\text{esquilo}}
\;+\;
\underbrace{\text{dissipação}}_{\text{gato}}
\;=\;
\text{costura}.
\end{equation}
Esta é \textbf{interpretação estrutural do projeto}:
a equação de NS não é ``definida'' pela nomenclatura
gato/esquilo na literatura clássica.

A associação
\[
\text{transporte} \leftrightarrow \text{esquilo},\qquad
\text{dissipação} \leftrightarrow \text{gato}
\]
é uma \textbf{escolha interpretativa da linguagem do projeto}
e não uma identificação entre as propriedades energéticas
esses termos e as definições internas de gato/esquilo.
O transporte incompressível satisfaz a identidade de
cancelamento energético sob as hipóteses usuais;
a viscosidade é o termo dissipativo.

%==========================================================================
\section{Identidade de energia}\label{sec:energia}
%==========================================================================

A identidade de energia~(\ref{eq:energia-padrao}) é estimativa fundamental,
mas \textbf{isoladamente não fecha a regularidade global 3D}.
A estimativa de energia em $L^2$, embora fundamental,
não é suficiente por si só para concluir existência global
e regularidade suave em três dimensões.

\subsection{Identidade discreta (MEDIDO NO PROJETO)}
No anel $\mathbb Z/P\mathbb Z$ com $P=65537$, $M=16$ modos,
Laplaciano $\Delta u_i = u_{i+1}-2u_i+u_{i-1}$ e
transporte em forma fluxo:
\begin{equation}\label{eq:energia-discreta}
\sum_i u_i\Delta u_i
=
-\sum_i(u_{i+1}-u_i)^2,\qquad
\sum_i\text{fluxo}(u_i^2)=0.
\end{equation}
Isto é \textbf{exato no modelo discreto testado}
(medido em \code{tests/navier_corpo.js}, $8{:}0$).
\textbf{Não é prova de regularidade do PDE contínuo.}

\medido{\code{tests/navier_corpo.js} ($8{:}0$; assinado por $M$).}

%==========================================================================
\section{Transformada Universal}\label{sec:transformada}
%==========================================================================

A Transformada Universal do projeto é (transformada.tex~\S~575):
\begin{equation}\label{eq:transformada}
\mathcal{G} = \mathcal{F}\circ\Lambda,\qquad
(\Lambda x)(u)=x(e^u),
\end{equation}
com $\mathcal{F}$ Fourier no eixo aditivo (fase) e
$\Lambda$ a mudança $x=e^u$ que converte escala
multiplicativa em soma aditiva. No eixo multiplicativo
corresponde a Mellin.

Papel no paper:
\begin{equation}\label{eq:transformada-papel}
\text{campo/representação}
\;\xrightarrow{\;\mathcal{G}\;}\;
\text{estrutura espectral}.
\end{equation}
A Transformada \textbf{não prova} regularidade: é ferramenta
de representação/análise dentro da arquitetura.

%==========================================================================
\section{Modelo discreto do projeto}\label{sec:discreto}
%==========================================================================

Além da identidade de energia, os testes do projeto
documentam:
\begin{itemize}
\item \code{tests/histerese_navier.js} ($6{:}0$): passo de
  calor como retencao espectral, espectro $\{1,\tfrac12,0\}$
  no ciclo $M=4$;
\item \code{tests/batuta_continuo.js} ($6{:}0$): segmento
  afim $B(s,t)=p(t)+s\,u(t)$; gap
  $(\Delta c)^2+(\Delta s)^2=2-t_1$ ($16/16$ no anel);
\item \code{tests/curvatura.c}: a borda $x^2=nx+1$ fecha sem
  tolerância --- o gap de Yang--Mills por andar é o
  comprimento da batuta entre ticks.
\end{itemize}
Estes resultados são \textbf{medidas do projeto}: testam
identidades internas no modelo discreto.
\textbf{Não constituem validação externa da questão do Clay.}

\medido{\code{tests/histerese_navier.js} ($6{:}0$);
\code{tests/batuta_continuo.js} ($6{:}0$);
\code{tests/curvatura.c}.}

%==========================================================================
\section{O resultado anunciado pela OpenAI (set.~2026)}
\label{sec:openai}
%==========================================================================

Segundo o anúncio oficial da OpenAI (8 set.~2026),
a OpenAI apresentou um resultado que afirma resolver
as partes C e D da formulação oficial do problema
de Navier--Stokes (existência e suavidade). A fonte
primária é o anúncio oficial de 8 de setembro de
2026 da OpenAI.

NÃO afirmamos que este resultado seja consenso,
teorema demonstrado ou validação independente. A
classificação aqui é apenas: \textbf{anunciado}.

\subsection{O que é afirmado (anunciado)}
\begin{enumerate}
\item \textbf{Afirmativa}: existe tempo finito $T$ com
  $\max|u(t)|\to\infty$ quando $t\to T^{-}$, energia
  cinética $\tfrac12\|u\|^2$ limitada.
\item \textbf{Formalização}: Lean; verificação adicional
  reportada.
\end{enumerate}

\subsection{Trabalho contemporâneo de Alpöge/Buckmaster}
Buckmaster e Alpöge publicaram preprints sobre
blowup forçado para Euler 3D (comentário externo;
sem entrada bibliográfica verificável neste paper).
Este trabalho é distinto do anúncio da OpenAI e
trata de Euler com força externa, não de Navier--Stokes
sem forçante. O estatuto de prioridade é documentado,
não decidido aqui.

\subsection{O que permanece independente}
\begin{itemize}
\item Um \textbf{anúncio} não é consenso matemático independente.
  Revisão por pares, críticas e lacunas são etapas separadas.
\item O anúncio da OpenAI afirma tratar as partes C e D
  da formulação oficial do problema. Seu estatuto
  matemático independente permanece uma questão
  distinta do anúncio.
\item Tao classificou o resultado de Euler como ``notável'';
  comentário externo, não validação do resultado da OpenAI.
\end{itemize}

\subsection{Classificação}
Tratamos o resultado como \textbf{CLAIMED} (anunciado,
não como teorema consolidado): a menos que fontes
independentes confirmem, mantemos \emph{anunciado/claimed}.

%==========================================================================
\section{Leitura do projeto sobre o resultado anunciado}
\label{sec:leitura-openai}
%==========================================================================

Na leitura interna do projeto, o fenómeno anunciado
interessa pela competição:
\[
\text{transporte/concentração}
\quad\leftrightarrow\quad
\text{dissipação}.
\]

Na leitura interna do projeto, o fenómeno anunciado é
interpretado heuristicamente como competição entre
concentração (esquilo) e dissipação (gato): o esquilo
concentra energia num núcleo que se contrai; o gato
dissipa no bordo. Na construção anunciada, a velocidade
diverge enquanto a energia permanece finita --- o gato
não impede a divergência, e o esquilo domina a escala
interna do núcleo.

\textbf{Não afirmamos} que esta decomposição seja
reformulação formal da prova da OpenAI, nem que a prova
confirme o Corpo Universal. A leitura é interpretação
do projeto.

%==========================================================================
\section{Costura}\label{sec:costura}
%==========================================================================

Chamamos de ``costura'' a leitura estrutural pela qual os
componentes conservativo e dissipativo são observados
conjunta\-mente no mesmo campo.

% Nota: \S{thm:costura} do projeto demonstra o seguinte:
% três domínios (mecânica, Carnot, Hopfield) partilham o
% mesmo invariante $Q$; o espelho $T$ vive nos três;
% o mesmo empurrão (translação / gérmen aditivo da fronteira)
% quebra os três. É teorema do projeto, não definição
% clássica de Navier--Stokes.

A identidade de energia~(\ref{eq:energia-discreta}) no
modelo discreto é a evidência computacional desta costura;
a sua extensão ao contínuo 3D é \textbf{declarada como
leitura}, não como teorema demonstrado.

%==========================================================================
\section{Fronteira e histerese}\label{sec:fronteira}
%==========================================================================

São conceitos canónicos internos do projeto:
\begin{description}
\item[fronteira] conceito formal interno da arquitetura do projeto:
  ``onde o irracional acaba''; é teorema interno, não lacuna.
\item[histerese] memória da borda / dependência do caminho;
  ``dobra $+$ seletor sobre o tick''.
\end{description}

\textbf{Não} afirmamos automaticamente que ``a fronteira é
exatamente a fronteira do problema de Navier--Stokes de Clay''.
São afirmações de níveis distintos.

Neste paper, investigamos se esta linguagem fornece uma
leitura útil da passagem do modelo discreto para o
regime contínuo.

\medido{\code{tests/histerese_navier.js} ($6{:}0$).}

%==========================================================================
\section{Vórtice e geometria}\label{sec:vortice}
%==========================================================================

A construção anunciada descreve um núcleo vortical que se
contrai e se alonga. O modelo do projeto utiliza uma
representação toroidal para determinados regimes vorticais;
esta identificação é \textbf{construção geométrica do modelo}
e não teorema topológico geral sobre vórtices de Navier--Stokes.

%==========================================================================
\section{Correspondência com Yang--Mills}\label{sec:yang-mills}
%==========================================================================

A conexão entre Navier--Stokes e Yang--Mills é uma
\textbf{correspondência estrutural interna / hipótese
estrutural do projeto} (Lei~2 $K^{**}=K$; Lei~4
$T+T^{\dg}$, $|\det|=1$).

\textbf{Não afirmamos} que o resultado da OpenAI resolva,
reduza ou prove qualquer parte do problema Yang--Mills
do Clay. Se não houver ponte matemática demonstrada nos
arquivos do projeto, não fabricamos uma.

%==========================================================================
\section{Tabela de classificação}\label{sec:tabela}
%==========================================================================

\begin{center}\small
\begin{tabular}{p{0.22\textwidth}p{0.42\textwidth}p{0.24\textwidth}}
\toprule
CAMADA & EXEMPLO & STATUS \\
\midrule
Matemática padrão & identidade de energia NS,
  $\partial_t u+(u\cdot\nabla)u=\nu\Delta u-\nabla p+f$ &
  literatura estabelecida \\
Resultado do modelo discreto &
  $\sum u\Delta u=-\sum(\Delta u)^2$, exato em $\mathbb Z/P\mathbb Z$ &
  MEDIDO NO PROJETO \\
Definição canônica do projeto &
  $T+T^{\dg}$, $|\det|=1$, gato/esquilo, costura &
  LEITURA DO PROJETO \\
Interpretação &
  gato/esquilo aplicado a NS, leitura espectral &
  interpretação estrutural \\
Hipótese/conjectura &
  fronteira contínua 3D, regularidade global &
  FRONTEIRA/CONJECTURA \\
Construção geométrica do projeto &
  toro como modelo de vórtice &
  modelo interpretativo \\
Ponte interna &
  correspondência Yang--Mills &
  hipótese do projeto \\
Resultado externo anunciado &
  OpenAI set.~2026 &
  CLAIMED; sujeito a fonte e validação \\
\bottomrule
\end{tabular}
\end{center}

%==========================================================================
\section{Estaduto}\label{sec:estaduto}
%==========================================================================

\begin{enumerate}
\item No vocabulário do projeto, a formulação cristalizada
representa uma tensão entre as componentes
conservativa e dissipativa. Esta é uma leitura
interna da arquitetura e não uma afirmação de
má colocação no sentido matemático usual.
\item Na linguagem interna do projeto, a formulação
  é organizada pela Lei~4 da tetral; nessa
  formulação a metade conservativa
  \textbf{fecha} (Liouville, $|\det|=1$ exacto, medido).
\item A regularidade global 3D é a \textbf{borda aberta} do
  Clay --- na linguagem interna do projeto, a passagem
  ao contínuo tridimensional é lida como uma fronteira
  da arquitetura, associada ao conceito interno de
  histerese; o projeto aponta a fronteira, não a fecha
  (não alegada).
\item O resultado anunciado pela OpenAI é uma
  \textbf{afirmação externa} sobre a formulação
  oficial do Clay (partes C e D); não prova das
  formulações não-forçadas (A/B).
\item Os testes discretos do projeto verificam identidades
  internas; não constituem validação independente do PDE
  contínuo.
\end{enumerate}

%==========================================================================
\section{Conclusão}\label{sec:conclusao}
%==========================================================================

Para a arquitetura do projeto, o anúncio fornece um
novo objeto externo sobre o qual a leitura de Lei~4/
costura pode ser aplicada; ele não transforma essa
leitura em prova do próprio formalismo.

Antes da construção, a pergunta era ``pode a dinâmica 3D
desenvolver uma singularidade?''. Depois do anúncio, a
pergunta passa a ser: o mecanismo, o enunciado exacto e a
formalização resistem ao escrutínio independente?

Para a arquitetura do Corpo Universal, a leitura permanece:
\begin{equation}\label{eq:costura-final}
\text{Lei~4}
\longrightarrow
\text{transporte}
+\text{dissipação}
\longrightarrow
\text{costura}
\longrightarrow
\text{fronteira}.
\end{equation}

\begin{quote}
\textbf{O projeto não reivindica que sua linguagem substitua
a teoria clássica de Navier--Stokes. Registra uma leitura da
equação de campo através da Lei~4, da Transformada e da
estrutura de costura, apoiada por identidades exactas no
modelo discreto e separada explicitamente das questões ainda
abertas no contínuo tridimensional.}
\end{quote}

A disputa de prioridade (Buckmaster/Alpöge vs OpenAI), o
estado da revisão por pares e a distinção A/B vs C/D são
etapas externas a este paper: este paper registra a
leitura interna, não decreta a resolução do Clay.

\medido{OpenAI (set. 2026): resultado anunciado.
Leitura Corpo Universal/Arquitetura: Lei~4,
gato/esquilo, histerese, fronteira interna.}

\begin{thebibliography}{9}
\bibitem{openai2026}
OpenAI (set.~2026),
\textit{An OpenAI model proposes a solution to the Navier--Stokes problem},
anúncio oficial de 8 de set.~2026.
\end{thebibliography}
\end{document}